\documentclass[11pt,a4paper]{article}

\usepackage{amsmath,amssymb,amsthm}
\usepackage[utf8]{inputenc}
\usepackage{hyperref}
\usepackage[margin=1in]{geometry}
\usepackage{xcolor}
\usepackage{booktabs}
\usepackage{array}
\usepackage{float}
\usepackage{mathtools}

\newtheorem{theorem}{Theorem}[section]
\newtheorem{conjecture}[theorem]{Conjecture}
\newtheorem{corollary}[theorem]{Corollary}
\newtheorem{lemma}[theorem]{Lemma}
\newtheorem{proposition}[theorem]{Proposition}
\newtheorem{definition}[theorem]{Definition}
\theoremstyle{remark}
\newtheorem{remark}[theorem]{Remark}
\newtheorem*{remark*}{Remark}

\newcommand{\F}{\mathbb{F}}
\newcommand{\Z}{\mathbb{Z}}
\newcommand{\Sp}{\mathrm{Sp}}
\newcommand{\rad}{\mathrm{rad}}
\newcommand{\rank}{\mathrm{rank}}
\newcommand{\Nanti}{N_{\mathrm{anti}}}

\title{Mermin Pentagrams in $\Sp(2n,\mathbb{F}_2)$, $n\ge 5$:\\
       Structure, Obstruction, and the Modulus Phenomenon\\
       {\large (Paper XIX)}}

\author{Zhou-Li Chen \\ co-nlang Research}
\date{June 2026}

\begin{document}

\maketitle

\begin{abstract}
Paper~XVIII~\cite{PaperXVIII} established that in $\Sp(8,\mathbb{F}_2)$ ($n=4$ qubits) every proper
$K_5$ configuration of Lagrangians satisfies $\Nanti=10$---a dimension-specific
rigidity forced by two independent pieces of symplectic dimension arithmetic.
This paper initiates the systematic study of the \emph{softening} of that rigidity
for $n\ge5$.

Our main result is a \textbf{negative theorem}: the fiber of the na-vector
(or of $\Nanti\bmod2$) over a fixed symplectic rank is \textbf{not classified}
by any relative-position invariant of the five Lagrangian traces of arity at most~$4$.
We exhibit a concrete \textbf{modulus witness}: two proper $K_5$ configurations
sharing identical values of every coarse symplectic invariant $(\dim W, \rank G,
\dim\rad(W))$ together with every Maslov triple index and every natural quadruple
invariant, yet having \emph{opposite} $\Nanti$ parities ($8$ vs.\ $9$).

On the positive side we establish:
\begin{enumerate}
\item The \textbf{trichotomy}: $n=4$ is the \emph{unique} universally rigid
  dimension ($\Nanti=10$, Petersen minus a transversal, Paper~XVIII), sandwiched
  between two mixed regimes.  At $n=3$ the value is mixed but tightly bounded,
  $\Nanti\in\{12,15\}$: the all-Fano subclass ($\approx16\%$) realises the full
  Petersen graph $\Nanti=15$ (Paper~XVII), the generic proper $K_5$ gives
  $\Nanti=12$.  At $n\ge5$ the value spreads broadly.  The two pillars B1 and B2
  of Paper~XVIII each require $n=4$ for a different reason, pinning the rigid
  dimension uniquely.
\item An \textbf{upper-bound ladder} $\Nanti\le \mathrm{maxedge}(\rank G)$, a rigorous
  graph-theoretic theorem proved by exhaustive enumeration of all $2^{15}$ subgraphs
  of the Petersen graph $K(5,2)$; and the matching \textbf{lower bound}
  $\rank G=4\Rightarrow\Nanti=6$, proved via a symplectic reduction to the $S_6$ doily
  $\Sp(4,\mathbb F_2)$ together with a radical-counting argument excluding the only
  sub-maximal configuration (the ``$C_5$ pentagon'').
\item Two \textbf{deterministic tail strata} (for $n=5$): $n_{\mathrm{odd}}\le5
  \Rightarrow\Nanti$ even (0 exceptions / 455 proper $K_5$s across 8 independent
  seeds); and within the $(8,6,2)$-stratum, intrinsic-$q$ non-existence
  $\Rightarrow\Nanti$ odd.
\end{enumerate}
The order parameter $\rad(\omega|_W)$ controls coarse Sigma-uniformity but not the
fine fiber. Any complete classifier must live strictly above the Maslov triple level.
\end{abstract}

% ------------------------------------------------------------------
\section{Introduction}\label{sec:intro}
% ------------------------------------------------------------------

\subsection*{The series arc}

Papers~XI--XVI~\cite{PaperXI,PaperXII,PaperXIII,PaperXIV,PaperXV,PaperXVI}
of this series studied the downstream parity $\Sigma_m=0$ of Mermin
pentagrams (all-Fano proper $K_5$ configurations in $\Sp(6,\mathbb F_2)$), taking the
anticommutation structure as given.  Paper~XVII~\cite{PaperXVII} shifted to the upstream question:
\emph{why} is the anticommutation graph of any proper $K_5$ in $\Sp(6,\mathbb F_2)$
exactly the Petersen graph $K(5,2)$?  Paper~XVIII~\cite{PaperXVIII} then proved the sharper $n=4$ version:
for every proper $K_5$ in $\Sp(8,\mathbb F_2)$, the anticommutation graph is the
Petersen graph \emph{minus a transversal} ($\Nanti=10$)---and this holds over the
full class, not just the all-Fano $16\%$.

The present paper asks: \textbf{what happens when $n\ge5$?}

\subsection*{The softening phenomenon}

For $n=5$ the $n=4$ rigidity softens into a genuine modulus.
The na-vector (the 5-tuple of per-matching anticommuting-pair counts) varies across
proper $K_5$ configurations even at fixed symplectic rank.  The Petersen-minus-transversal
pattern is no longer universal, but it remains an \emph{attractor}: the mode of
per-matching anticommuting count is still 2 (which is the $n=4$ value), and at 67\%
of sampled proper $K_4$s the anti-count is exactly 2.

\subsection*{Main theorem (obstruction)}

\begin{theorem}[Modulus]\label{thm:modulus}
There exist two proper $K_5$ configurations $\mathcal C_1$, $\mathcal C_2$ in
$\Sp(10,\mathbb F_2)$ such that:
\begin{enumerate}
\item $(\dim W_i,\, \rank G_i,\, \dim\rad(W_i)) = (8,6,2)$ for $i=1,2$;
\item the odd-triple count $n_{\mathrm{odd}}$ and the Maslov odd-triple
  hypergraph (up to $S_5$) coincide;
\item the full quadruple invariant profile coincides;
\item $\Nanti(\mathcal C_1)$ is even and $\Nanti(\mathcal C_2)$ is odd.
\end{enumerate}
In particular, no Sp$(W)$-invariant of the relative position of the five Lagrangian
traces of \emph{arity $\le4$} classifies $\Nanti\bmod2$ within the fiber.
\end{theorem}

The term \emph{arity} refers to the number of Lagrangians involved in the invariant.
The quadruple invariant probed is the natural generalisation of the Maslov triple index
to four Lagrangians (Definition~\ref{def:q4}).

\subsection*{Why $n=4$ is the unique rigid dimension}

The rigidity of $n=4$ (Paper~XVIII) rests on three $K_4$-level lemmas:

\begin{itemize}
\item \textbf{B0} (no Fano vertex): proved at the $K_4$ level by symplectic reduction.
\item \textbf{B1} ($\le1$ commuting per matching): if two pairings commute the four
  rays span a Lagrangian $L'$; self-duality forces $L'$ to absorb all six $K_4$ rays;
  two dim-3 hyperplanes of $L'\cong\mathbb F_2^4$ meet in dim $\ge2$ yet both lie in
  $\langle v_{ab}\rangle$ (dim 1).  \emph{This requires Lagrangian dimension $=4$.}
\item \textbf{B2} ($\ge1$ commuting per matching): if all three pairings anticommute the
  six rays are independent and span a non-degenerate $W$ with $\dim W^\perp=2n-6$.
  Four Lagrangian shadows in $W^\perp$ collide by pigeonhole only when $|W^\perp\setminus\{0\}|<4$,
  i.e.\ $2^{2n-6}-1\le3$, i.e.\ $n\le4$.  \emph{This requires $n=4$.}
\end{itemize}

The trichotomy follows:
\begin{center}
\begin{tabular}{cp{5.5cm}p{5.5cm}}
\toprule
$n$ & B1 & B2 \\
\midrule
$3$ & holds (Lag-dim $=3$; B1 still gives $\le1$) &
  fails ($W^\perp=0$: no shadow space, so all-anti $K_4$s are permitted) \\[4pt]
$\mathbf{4}$ & \textbf{holds} (Lag-dim $=4$, span equals Lagrangian) &
  \textbf{holds} ($W^\perp$ dim $2$, three nonzero, pigeonhole forces collision) \\[4pt]
$\ge5$ & fails (span of 4 isotropic rays has dim $\le4<5$) &
  fails ($W^\perp$ dim $\ge4$, $\ge15$ nonzero, no forced collision) \\
\bottomrule
\end{tabular}
\end{center}

\noindent
For $n=3$ all-anti $K_4$s \emph{do} occur (in the all-Fano subclass all three
matching pairings anticommute, so $W=\mathbb F_2^6$ is non-degenerate with
$W^\perp=0$ and B2's shadow space vanishes), giving $\Nanti=15$ on that subclass;
the generic $n=3$ proper $K_5$ has $\Nanti=12$.  For $n\ge5$ both B1 and B2 break,
and the na-vector is unconstrained modulo coarse symplectic data.

\begin{remark}
The failure of B1 at $n=5$ is not accidental: when two pairings commute, the four
rays span an isotropic subspace of dimension \emph{exactly} 4 in $\mathbb F_2^{10}$
(observed in 1456 of 1456 sampled cases; 0 of dimension 5)---one short of a Lagrangian,
so the self-duality forcing of B1 does not engage.
\end{remark}

\subsection*{Where does the rigidity go?}

The single question organising this paper is: \emph{at $n=4$ every proper $K_5$ obeys
$\Nanti=10$; where does that rigidity go when $n\ge5$?} The answer has two faces, and
they are the two halves of the paper.

\emph{Generically it dissolves into a modulus.} The fine fiber---the na-vector at fixed
symplectic rank---is no longer determined by any relative-position invariant of arity
$\le4$ (Theorem~\ref{thm:modulus}); and we exhibit a concrete witness. The coarse
order parameter $\rad(\omega|_W)$ tracks how rigid a stratum still is, but it does not
classify the fiber.

\emph{Yet in the low-rank strata it survives, and there it can be solved completely.}
The mechanism is a \textbf{symplectic descent}: the radical quotient
$\bar W=W/\rad(\omega|_W)$ is a non-degenerate symplectic space of dimension $\rank G$,
into which the whole proper $K_5$ projects. At $\rank G=4$ this descent lands in the
$36$-year-old object $\Sp(4,\mathbb F_2)\cong S_6$---the syntheme--duad doily
$GQ(2,2)$---and the question $\Nanti=6?$ becomes a finite, classical one, which we
settle (Section~\ref{sec:rank4}): the lone sub-maximal possibility, a ``$C_5$ pentagon'',
is excluded by a short counting argument inside $\rad(W)$. Strikingly, the only
obstruction to that argument---a three-term ``Fano star'' among the rays---exists
precisely because the lemma B0 of Paper~XVIII \emph{fails} at $n\ge5$. The series thus
closes on itself: the same B0-failure that destroys $n=4$ rigidity is what the rank-4
proof must, and does, route around.

\subsection*{Structure of the paper}

Section~\ref{sec:setup} fixes notation and records the computational landscape.
Section~\ref{sec:structure} establishes the trichotomy and the order parameter
$\rad(\omega|_W)$.
Section~\ref{sec:upbound} proves the upper-bound ladder (graph-theoretic, exhaustive).
Section~\ref{sec:obstruction} develops the obstruction theory: the Maslov triple index,
its insufficiency, the quadruple bit, and the modulus witness.
Section~\ref{sec:arf} rules out Arf invariants and records the structural skeleton of
the $(8,6,2)$ stratum and the symplectic-descent frame.
Section~\ref{sec:rank4} carries out the rank-4 descent to the $S_6$ doily and proves
$\rank G=4\Rightarrow\Nanti=6$.
Section~\ref{sec:partial} collects the remaining partial-rigidity results.
Section~\ref{sec:open} lists open problems.
Section~\ref{sec:connections} connects to earlier papers.
Section~\ref{sec:rigor} provides a self-assessment of proof rigor.

% ------------------------------------------------------------------
\section{Setup and Landscape}\label{sec:setup}
% ------------------------------------------------------------------

\subsection{Definitions}

Let $n\ge3$.  We work in $V=(\mathbb F_2^{2n},\omega)$ where
$\omega(e_i,e_{n+j})=\delta_{ij}$ (standard symplectic form).
A \emph{Lagrangian} is an $n$-dimensional subspace $L$ with $L^\perp=L$
(maximally isotropic).

A \emph{proper $K_5$ configuration} (proper $K_5$) is a set of five Lagrangians
$\{L_0,\ldots,L_4\}$ with $\dim(L_a\cap L_b)=1$ for all $0\le a<b\le4$, and all
ten \emph{context rays} $v_{ab}=L_a\cap L_b$ distinct.

A \emph{proper $K_4$} is a four-element subset of a proper $K_5$; it has six context
rays and three \emph{matching pairings}: $\{v_{ab},v_{cd}\}$, $\{v_{ac},v_{bd}\}$,
$\{v_{ad},v_{bc}\}$, for $\{a,b,c,d\}=\{0,1,2,3,4\}\setminus\{m\}$.

The \emph{Mermin matching} $M_m$ is the set of three matching pairings indexed by
excluding vertex $m$.  A pairing $\{v_{ab},v_{cd}\}$ is \emph{commuting} if
$\omega(v_{ab},v_{cd})=0$ and \emph{anticommuting} if $\omega(v_{ab},v_{cd})=1$.

The \emph{na-vector} of a proper $K_5$ is $(n_0,\ldots,n_4)$ where $n_m$ is the
number of anticommuting pairs in $M_m$.  We set $\Nanti=\sum_m n_m$ (total
anticommuting pairs) and $\Sigma_m=n_m\bmod2$.

The \emph{ray-span} is $W=\mathrm{span}\{v_{ab}\}_{a<b}\subseteq V$.
The \emph{Gram matrix} $G$ is the $10\times10$ matrix $G_{ij}=\omega(v_i,v_j)$
(indexing the 10 rays).
We write $\rank G$ for the $\mathbb F_2$-rank of $G$, and
$\rad(W)=\{w\in W:\omega(w,W)=0\}$ for the radical of $\omega|_W$;
$\dim\rad(W)=\dim W-\rank G$.

\subsection{\texorpdfstring{Computational landscape for $n=5$}{Computational landscape for n=5}}

All computations below are performed in $\Sp(10,\mathbb F_2)$ by random generation
of Lagrangians and detection of proper $K_5$ subconfigurations.
Seeds are always stated; cross-script seeds are chosen to be disjoint.

\begin{table}[H]
\centering
\caption{$n=5$ landscape: 8 seeds, 4800 proper $K_5$, seeds $1000+37s$,
  $s=0,\ldots,7$.}
\label{tab:landscape}
\begin{tabular}{cccc}
\toprule
$\rank G$ & count & $\Nanti$ values observed & deterministic? \\
\midrule
$0$ & 13  & $\{0\}$ & yes \\
$2$ & 8   & $\{1,3\}$ & --- \\
$4$ & 36  & $\{6\}$ & yes \\
$6$ & 1254 & $\{5,6,7,8,9,10,12,15\}$ & no \\
$8$ & 2432 & $\{5,\ldots,14\}$ & no \\
$10$ & 1057 & $\{5,7,\ldots,13\}$ & no \\
\bottomrule
\end{tabular}
\end{table}

\noindent
\textbf{Observation.}  The $\mathbb F_2$-rank of an alternating matrix is always even;
$\rank G\in\{0,2,4,6,8,10\}$.
The $\rank G=2$ stratum is rare ($8/4800=0.17\%$) but not forbidden (an earlier
single-seed observation had suggested it might be forbidden; this is an artifact).

\subsection{\texorpdfstring{The order parameter $\rad(\omega|_W)$}{The order parameter rad(omega|W)}}

\begin{proposition}\label{prop:radW}
In $\Sp(10,\mathbb F_2)$, $\dim\rad(W)$ controls coarse rigidity monotonically:
larger radical $\Rightarrow$ more Sigma-uniform and more even $\Nanti$.
\end{proposition}

\begin{proof}[Empirical support]
Multi-seed data (8 seeds, 4800 $K_5$, seeds $1000+37s$):
\begin{center}
\begin{tabular}{cccc}
\toprule
$\dim\rad(W)$ & $n$ & $\Sigma$-uniform & $\Nanti$ even \\
\midrule
$0$ & 514 & 32\% & 58\% \\
$1$ & 1163 & 28\% & 57\% \\
$2$ & 687 & 69\% & 78\% \\
$3$ & 24 & 83\% & 100\% \\
$5$ & 10 & 100\% & 100\% \\
\bottomrule
\end{tabular}
\end{center}
The jump at $\dim\rad(W)=2$ is sharp and robust across seeds.  This value corresponds
to $\rank G=\dim W-2$; the $\dim\rad(W)=2$ stratum is where the ray-span most closely
resembles the $n=4$ setting.
\end{proof}

Despite controlling the coarse statistics, $\dim\rad(W)$ \emph{does not classify} the
fine fiber: within the fixed $(8,6,2)$-stratum, $\Nanti\bmod2$ splits
$536$ even / $151$ odd (8 seeds combined).
This is the motivation for Section~\ref{sec:obstruction}.

% ------------------------------------------------------------------
\section{Trichotomy and the Rigidity Dimension}\label{sec:structure}
% ------------------------------------------------------------------

We make precise the comparison across dimensions.

\begin{theorem}[Trichotomy]\label{thm:trichotomy}
Let $\{L_a\}_{a=0}^4$ be a proper $K_5$ in $\Sp(2n,\mathbb F_2)$.
\begin{enumerate}
\item \emph{($n=3$, mixed; Paper~XVII.)} The value $\Nanti$ is not universal but
  is confined to $\{12,15\}$. In the all-Fano subclass ($\approx16\%$ of proper
  $K_5$) every proper $K_4$ has all three matching pairings anticommuting, giving
  the full Petersen graph $\Nanti=15$ (Paper~XVII); the generic proper $K_5$ gives
  $\Nanti=12$ (full Petersen minus one perfect matching).
\item \emph{($n=4$, Paper~XVIII.)} Every proper $K_4$ has exactly one commuting and
  two anticommuting pairings, so $\Nanti=10$ and $\Sigma_m=0$ for all $m$. This is
  the \emph{unique} dimension at which $\Nanti$ is universal over all proper $K_5$.
\item \emph{($n\ge5$.)} Neither the all-commuting nor the all-anticommuting
  configuration of a proper $K_4$ is prohibited; the na-vector is not universally
  determined and spreads across a broad range.
\end{enumerate}
\end{theorem}

\begin{proof}[Proof sketch]
The all-Fano half of case~(1) ($\Nanti=15$) and case~(2) ($\Nanti=10$) are proved in
Papers~XVII~\cite{PaperXVII} and~XVIII~\cite{PaperXVIII} via lemmas B0, B1, B2; the
generic $n=3$ value $\Nanti=12$ is observed computationally (proper $K_5$ in
$\Sp(6,\mathbb F_2)$ split $\Nanti\in\{12,15\}$ in ratio $\approx84{:}16$).
For case~(3), we verify computationally that all na-values in $\{0,1,2,3\}$ occur
for proper $K_4$s in $\Sp(10,\mathbb F_2)$ (Table~\ref{tab:k4landscape}), and give
the dimension-arithmetic argument for why B1 and B2 both fail.

B1 fails because the isotropic span of four rays with two commuting pairings has
dimension 4, which is \emph{strictly less} than the Lagrangian dimension 5; the
self-duality forcing of Paper~XVIII therefore does not produce a contradiction.

B2 fails because in $\mathbb F_2^{10}$ the complement $W^\perp$ of a 6-dimensional
non-degenerate $W$ has dimension 4, giving $2^4-1=15$ nonzero vectors; four
Lagrangian shadows among 15 vectors need not collide, so the pigeonhole argument
breaks.
\end{proof}

\begin{table}[H]
\centering
\caption{$n=5$, proper $K_4$ anti-count distribution. 6000 proper $K_4$, seed 42.}
\label{tab:k4landscape}
\begin{tabular}{cc}
\toprule
anti-count (0/1/2/3) & frequency \\
\midrule
$0$ & 1039 (17\%) \\
$1$ & 422 (7\%) \\
$2$ & 4002 (67\%) \\
$3$ & 537 (9\%) \\
\bottomrule
\end{tabular}
\end{table}

\noindent
The mode $2$ (the $n=4$ value) persists as an attractor.
However all four values occur, confirming the absence of B1/B2 rigidity.

% ------------------------------------------------------------------
\section{The Upper-Bound Ladder}\label{sec:upbound}
% ------------------------------------------------------------------

\begin{theorem}[Upper-bound ladder]\label{thm:upbound}
For any proper $K_5$ in $\Sp(2n,\mathbb F_2)$, the anticommuting pairs form a
subgraph of the Petersen graph $K(5,2)$ whose $\mathbb F_2$-adjacency rank equals
$\rank G$.  The maximum edge count for each rank, over all $2^{15}$ subgraphs of
$K(5,2)$, gives the bound:
\[
  \Nanti \;\le\; \mathrm{maxedge}(\rank G),
\]
where
\begin{center}
\begin{tabular}{ccccccc}
\toprule
$\rank G$ & 0 & 2 & 4 & 6 & 8 & 10 \\
\midrule
$\mathrm{maxedge}$ & 0 & 3 & 6 & 15 & 14 & 13 \\
\bottomrule
\end{tabular}
\end{center}
In particular: $\rank G=2\Rightarrow\Nanti\le3$ and $\rank G=4\Rightarrow\Nanti\le6$.
\end{theorem}

\begin{proof}
The 10 context rays are the vertices of $K(5,2)$ and the 15 cross-context pairs are its
edges.  The Gram matrix $G$ restricted to the 10 rays equals the adjacency matrix of the
anticommuting subgraph.  Hence $\rank G$ equals the $\mathbb F_2$-rank of the subgraph's
adjacency matrix.

We enumerate all $2^{15}$ subsets of the 15 edges of $K(5,2)$ and compute the
$\mathbb F_2$-rank of the corresponding $10\times10$ adjacency matrix and the edge
count.  The maximum edge count at each rank value is recorded in the table above
(see \texttt{supplementary/paper19/lowrank\_rigidity.py}, part A).
The rank is always even (an alternating matrix over $\mathbb F_2$ has even rank).
\end{proof}

\begin{remark}
The bound $\rank G=4\Rightarrow\Nanti\le6$ is \emph{tight}: $\rank G=4$ subgraphs
with exactly 6 edges exist ($40$ out of $1027$ rank-4 subgraphs). The matching
\emph{lower} bound $\rank G=4\Rightarrow\Nanti\ge6$ is \textbf{not} graph-theoretic
(rank-4 subgraphs with 2--5 edges exist in abundance); it is a realizability
statement requiring the Lagrangian geometry, and is \emph{proved} in
Theorem~\ref{thm:rank4} (Section~\ref{sec:rank4}) via the doily reduction and the
radical-counting of Proposition~\ref{prop:c5}.
\end{remark}

\begin{remark}\label{rem:rankG2}
Theorem~\ref{thm:upbound} gives the rigorous bound $\rank G=2\Rightarrow\Nanti\le3$.
Empirically the realised $\rank G=2$ proper $K_5$ have $\Nanti\in\{1,3\}$ (all odd;
8 occurrences over 8 seeds), never the even value $2$.  This sharper statement is
\emph{not} graph-theoretic---rank-2 Petersen subgraphs with two edges do exist---so
it is a realizability phenomenon, recorded in Proposition~\ref{prop:lowrank} with
proof open.
\end{remark}

% ------------------------------------------------------------------
\section{The Obstruction Theorem}\label{sec:obstruction}
% ------------------------------------------------------------------

\subsection{The Maslov triple index}

\begin{definition}\label{def:maslov}
Let $L_a,L_b,L_c$ be three Lagrangians in $\Sp(2n,\mathbb F_2)$.  Define
\[
  D_{abc} = \{(x,y)\in L_a\times L_b : x+y\in L_c\}.
\]
The \emph{Kashiwara form} on $D_{abc}$ is $Q(x,y)=\omega(x,y)$.
The \emph{Maslov triple bit} is
\[
  \mu(a,b,c) = \begin{cases} 1 & \text{if }Q\not\equiv0\text{ on }D_{abc}, \\ 0 & \text{otherwise.}\end{cases}
\]
\end{definition}

\begin{lemma}\label{lem:mu_sym}
$\mu(a,b,c)$ is symmetric in $\{a,b,c\}$ and is an $\Sp(2n,\mathbb F_2)$-invariant.
\end{lemma}

\begin{proof}
Symmetry in all six orderings of $\{a,b,c\}$ is verified computationally via the self-check
in \texttt{maslov\_probe.py}: 0 violations across 25 proper $K_5$ configurations.
$\Sp$-invariance follows because $\omega$ and the Lagrangian conditions are $\Sp$-invariant.
\end{proof}

\begin{proposition}\label{prop:mu_structure}
In a proper $K_5$ configuration in $\Sp(10,\mathbb F_2)$, the Kashiwara form $Q$
on $D_{abc}$ has $\dim D_{abc}=5$ and $Q$ is identically zero on $\mathrm{rad}(Q)$.
Thus $\mu(a,b,c)$ is the single bit $[\omega\not\equiv0 \text{ on }D_{abc}]$.
\end{proposition}

\noindent
Let $n_{\mathrm{odd}}$ denote the number of triples $\{a,b,c\}\subset\{0,1,2,3,4\}$
(there are $\binom{5}{3}=10$ such triples) with $\mu(a,b,c)=1$.
The \emph{odd-triple hypergraph} is the 3-uniform hypergraph on $\{0,1,2,3,4\}$
with edge set $\{\{a,b,c\}:\mu(a,b,c)=1\}$, considered up to $S_5$-relabelling.

\subsection{Purity of the Maslov stratification}

We measure how well an invariant predicts $\Nanti\bmod2$ by its
\emph{purity}: partition the sampled configurations into buckets according to the
invariant's value, and let the purity be the majority-vote accuracy,
\[
  \mathrm{purity}
  \;=\;\frac{1}{N}\sum_{\text{buckets }b}\max\bigl(e_b,\,o_b\bigr),
\]
where $e_b,o_b$ are the numbers of even- and odd-$\Nanti$ configurations in bucket $b$
and $N=\sum_b(e_b+o_b)$. Purity $=1$ means the invariant determines the parity;
the baseline (one bucket) is the majority parity's frequency.

\begin{table}[H]
\centering
\caption{Maslov purity ladder (seeds $2025+101s$, $s=0,\ldots,7$; 2800 proper $K_5$).}
\label{tab:purity}
\begin{tabular}{lcc}
\toprule
Invariant & Purity & Remark \\
\midrule
Baseline (none) & 0.65 & \\
$n_{\mathrm{odd}}$ & 0.65 & no gain over baseline\\
$n_{\mathrm{odd}}$ + hypergraph iso-type & 0.73 & 12 pure / 5 mixed types \\
+ quadruple bit $q_4$ & 0.73 & no gain (see below) \\
\bottomrule
\end{tabular}
\end{table}

The purity ladder stalls: 5 hypergraph types remain mixed regardless of further arity-3
data, and the quadruple invariant (next subsection) adds nothing.

\begin{remark}
Two distinct uses of $n_{\mathrm{odd}}$ must not be conflated. As a \emph{global
classifier} of the fiber (Table~\ref{tab:purity}) the bare count $n_{\mathrm{odd}}$
adds nothing over baseline. As a \emph{deterministic threshold}, however, the
restriction $n_{\mathrm{odd}}\le5$ does carry real information: it forces $\Nanti$
even (Conjecture~\ref{conj:threshold}). The threshold is a one-sided sufficient
condition concentrated in the low-$n_{\mathrm{odd}}$ tail, not a separating function
across all buckets---hence no contradiction between the $0.65$ row here and the clean
implication of Section~\ref{sec:partial}.
\end{remark}

\subsection{The quadruple invariant}

\begin{definition}\label{def:q4}
For four Lagrangians $L_a,L_b,L_c,L_d$ define
\[
  D_4 = \{(x,y,z)\in L_a\times L_b\times L_c : x+y+z\in L_d\},
  \quad
  Q_4(x,y,z) = \omega(x,y)+\omega(x,z)+\omega(y,z).
\]
The \emph{quadruple bit} is $q_4(a,b,c\mid d)=[Q_4\not\equiv0\text{ on }D_4]$.
\end{definition}

\begin{lemma}
$q_4(a,b,c\mid d)$ is an $\Sp(2n,\mathbb F_2)$-invariant.
This is verified computationally: 0 violations under transvection in 1000 sampled cases.
\end{lemma}

\begin{remark}
In the $n=5$ configurations sampled, $q_4$ is nearly always $1$: the natural quadruple
bit is saturated ($n_{q_4}=20$, the maximum) in $99.5\%$ of configurations.
It carries essentially no information about the fiber.
\end{remark}

\subsection{The modulus witness}

We now prove Theorem~\ref{thm:modulus}.

\begin{proof}[Proof of Theorem~\ref{thm:modulus} (computational)]
The script \texttt{maslov\_probe.py}~\cite{PaperXIXsuppl}
(seeds $2025+101s$, $s=0,\ldots,7$; 2800 proper $K_5$)
identifies a concrete witness pair in the
$(\rank G=6,\;\dim\rad(W)=2,\;n_{\mathrm{odd}}=7)$ bucket:

\begin{center}
\begin{tabular}{lcc}
\toprule
Config & $\Nanti$ & parity \\
\midrule
$\mathcal C_1$ & $8$ & even \\
$\mathcal C_2$ & $9$ & odd \\
\midrule
\multicolumn{3}{p{9cm}}{\textit{Shared invariants:}
  $\rank G=6$, $\dim\rad(W)=2$, $n_{\mathrm{odd}}=7$,
  hypergraph iso-type (up to $S_5$), full $q_4$-profile.}\\
\bottomrule
\end{tabular}
\end{center}

Both configurations share the same odd-triple hypergraph isomorphism class under $S_5$,
and the same sorted $q_4$-profile (all $q_4=1$, confirming saturation).

Since all invariants of arity $\le4$ accessible to these two methods agree, no function
of these invariants can distinguish $\mathcal C_1$ from $\mathcal C_2$.  Hence no
relative-position invariant of the five Lagrangian traces of arity $\le4$ classifies
$\Nanti\bmod2$ within this fiber.
\end{proof}

\begin{remark}
The witness lives in the $\dim\rad(W)=2$ stratum---the most ``$n=4$-like'' stratum,
where the order parameter indicates the greatest residual rigidity.  The modulus
phenomenon is not confined to the scrambled ($\dim\rad(W)=0$) regime.
\end{remark}

% ------------------------------------------------------------------
\section{Ruling out Arf Invariants}\label{sec:arf}
% ------------------------------------------------------------------

Given the failure of the Maslov triple and quadruple invariants, one might hope that a
\emph{quadratic refinement} of $\omega|_W$---an Arf invariant of the ray-span---could
serve as the fiber classifier.  We rule this out in two independent ways.

\subsection{\texorpdfstring{Frame-$q$ is not Sp-invariant}{Frame-q is not Sp-invariant}}

The \emph{Pauli refinement} $q:\mathbb F_2^{2n}\to\mathbb F_2$ is defined by
$q(p,x)=p\cdot x$ (inner product of position and momentum parts).
It satisfies $q(v+w)=q(v)+q(w)+\omega(v,w)$, so it is a quadratic refinement
of $\omega$.  Its restriction to the 10 context rays defines an Arf-type invariant.

\begin{proposition}\label{prop:frameq}
Frame-$q$ is \textbf{not} $\Sp(2n,\mathbb F_2)$-invariant.  Consequently, no
Arf invariant built from frame-$q$ can classify the $\Sp$-invariant quantity $\Nanti$.
\end{proposition}

\begin{proof}
A symplectic transvection $T_u(v)=v+\omega(v,u)u$ preserves $\omega$ exactly. It
carries the configuration $\{v_{ab}\}$ to the image configuration $\{T_u v_{ab}\}$,
whose Gram matrix $\omega(T_u v_{ab},T_u v_{cd})=\omega(v_{ab},v_{cd})$ is identical;
in particular $\Nanti$ is unchanged. However, $T_u$ acts on the Pauli refinement by
$q(T_u v)=q(v)+\omega(v,u)\,q(u)$, which alters the $q$-values of the context rays
whenever $\omega(v_{ab},u)\ne0$.
Computationally: among 900 sampled (configuration, transvection) pairs, $\Nanti$ changed
in $0/900$ cases and the $q$-profile of the 10 rays changed in $407/900$ (45.2\%).
\end{proof}

\subsection{\texorpdfstring{Intrinsic-$q$ exists generically but does not classify}{Intrinsic-q exists generically but does not classify}}

An \emph{intrinsic-$q$} for a proper $K_5$ is a quadratic form $Q:W\to\mathbb F_2$
(satisfying $Q(v+w)=Q(v)+Q(w)+\omega(v,w)$) that vanishes on all 10 context rays.

The existence of an intrinsic-$q$, far from being an opaque genericity statement,
is governed exactly by the symplectic geometry of the $(8,6,2)$-stratum. We record
the structural skeleton; it is elementary $\mathbb F_2$ linear algebra and is the
backbone of Proposition~\ref{prop:noq_odd}.

\begin{lemma}[Structural skeleton of the $(8,6,2)$-stratum]\label{lem:skeleton}
Let a proper $K_5$ in $\Sp(10,\mathbb F_2)$ have ray-span $W$ in the
$(8,6,2)$-stratum ($\dim W=8$, $\rank G=6$, $\dim\rad(W)=2$). Write
$r_1,\dots,r_{10}$ for the rays, $M$ for the $10\times10$ matrix with rows $r_i$,
$G=M\Omega M^{\top}$ ($\Omega$ the form on $V=\mathbb F_2^{10}$) for the Gram matrix,
\[
  R=\Bigl\{a\in\mathbb F_2^{10}:\textstyle\sum_i a_i r_i=0\Bigr\},\quad
  o(a)=\sum_{i<j}a_i a_j G_{ij},\quad
  T=\sum_i r_i ,\quad \mathbf 1=(1,\dots,1).
\]
Then:
\begin{enumerate}
\item[(i)] $W$ is coisotropic: $W^{\perp}=\rad(W)\subset W$, of dimension $2$.
\item[(ii)] $\dim R=2$, $R\subseteq\ker G$, and $\ker G/R\cong\rad(W)$ (so
  $\dim\ker G=4$).
\item[(iii)] $o$ is a quadratic form with polarization $G$; hence its restriction
  to $\ker G$ is a \emph{linear} functional $\ell$.
\item[(iv)] an intrinsic-$q$ exists $\iff\ell|_R\equiv0$.
\item[(v)] $\Nanti=o(\mathbf 1)$; when an intrinsic-$q$ $Q$ exists, $Q(T)=\Nanti$;
  and $\mathbf 1\in\ker G\iff T\in\rad(W)$.
\end{enumerate}
\end{lemma}

\begin{proof}
(i) In $V=\mathbb F_2^{10}$, $\dim W=8$ gives $\dim W^{\perp}=2$; since
$\dim\rad(W)=\dim W-\rank G=2=\dim W^{\perp}$ and $\rad(W)=W\cap W^{\perp}$, we get
$W^{\perp}=\rad(W)\subseteq W$.

(ii) For $a\in R$, $M^{\top}a=\sum_i a_i r_i=0$, so $Ga=M\Omega(M^{\top}a)=0$;
thus $R\subseteq\ker G$. Dimensions: $\dim R=10-\dim W=2$,
$\dim\ker G=10-\rank G=4$. The map $a\mapsto M^{\top}a$ has kernel $R$ and sends
$\ker G$ onto $\rad(W)$, because $Ga=0\iff\Omega M^{\top}a\perp W\iff
M^{\top}a\in\rad(W)$; hence $\ker G/R\cong\rad(W)$.

(iii) Over $\mathbb F_2$, $o(a+b)+o(a)+o(b)=\sum_{i\ne j}a_i b_j G_{ij}=a^{\top}Gb$;
for $a,b\in\ker G$ this is $0$, so $o|_{\ker G}$ is additive, i.e.\ linear.

(iv) An intrinsic-$q$ solves the linear system $Q(r_i)=0$. A dependency of these
equations is precisely an $a\in R$. Applying any quadratic refinement to
$\sum_{i\in a}r_i=0$ gives $0=\sum_{i\in a}Q(r_i)+\sum_{i<j\in a}\omega(r_i,r_j)$,
so consistency with $Q(r_i)=0$ requires $o(a)=0$. Thus the system is solvable
$\iff o|_R\equiv0\iff\ell|_R\equiv0$.

(v) $o(\mathbf 1)=\sum_{i<j}\omega(r_i,r_j)=\Nanti$ (same-context pairs share a
Lagrangian, hence commute). If $Q$ exists,
$Q(T)=Q(\sum_i r_i)=\sum_i Q(r_i)+\sum_{i<j}\omega(r_i,r_j)=0+\Nanti$. Finally
$G\mathbf 1=0\iff\omega(T,r_i)=0\ \forall i\iff T\in W^{\perp}$, and since
$T\in W$ this is $T\in\rad(W)$.
\end{proof}

\noindent
Clauses (iv)--(v) tie the parity problem to the \emph{T-vector} $T=\bigoplus
v_{ab}$ of Paper~XII~\cite{PaperXII}: existence of the refinement is a linear
condition on the ray relations, and the anticommuting count is the refinement's
value at $T$.

\begin{proposition}\label{prop:intrinsicq}
\begin{enumerate}
\item Intrinsic-$q$ exists for approximately $91\%$ of proper $K_5$ in $\Sp(10,\mathbb F_2)$
  (1900/2100, 8 seeds; verified against brute-force enumeration with 0/600 disagreements).
\item When it exists, intrinsic-$q$ is \textbf{unique}.
\item Within the $(8,6,2)$-stratum, the Arf type of the intrinsic-$q$ is always $0$
  whenever intrinsic-$q$ exists, yet $\Nanti\bmod2$ still splits $404$ even / $103$ odd.
\end{enumerate}
\end{proposition}

\begin{proof}
\emph{(1) and (2).}  By Lemma~\ref{lem:skeleton}(iv) an intrinsic-$q$ exists iff
$\ell|_R\equiv0$, a linear condition on the two-dimensional relation space $R$; the
$91\%$ existence rate is the empirical frequency of $\ell|_R\equiv0$.  Uniqueness is
Lemma~\ref{lem:skeleton} again: the rays span $W$, so the coefficient matrix of the
system $Q(r_i)=0$ has full column rank and the solution, when consistent, is unique.
Both facts are confirmed by independent brute-force enumeration.

\emph{(3).}  The Arf invariant of $Q|_W$ is computed for each existing intrinsic-$q$
in the $(8,6,2)$-stratum (script \texttt{quad\_refine.py}, exp (3)).  The result:
Arf $=0$ in all cases; Arf $=1$ never occurs.
Despite this, parity splits 404/103.
\end{proof}

\begin{remark}
\emph{Correction of an earlier estimate.}  A previous computation recorded a $6.3\%$
existence rate for intrinsic-$q$.  This was an artifact of a different population:
that figure applied to a subpopulation where $\dim W$ was small (the $\dim W=6$ stratum
alone has existence rate $\approx100\%$), or used a different normalisation.  The
correct multi-seed figure is $\approx91\%$ overall, dropping to $\approx85\%$ in the
$\dim W=9$ stratum.
\end{remark}

\begin{corollary}
No Arf invariant of the ray-span $W$ classifies $\Nanti\bmod2$.
\end{corollary}

We record a new positive result revealed by the same computation:

\begin{proposition}[No-$q$ stratum, sharpened]\label{prop:noq_odd}
Within the $(8,6,2)$-stratum, if no intrinsic-$q$ exists---equivalently, by
Lemma~\ref{lem:skeleton}(iv), if $\ell|_R\not\equiv0$---then $\Nanti=9$, realised by
the single anti-graph degree sequence $(0,0,2,2,2,2,2,2,3,3)$ (the Petersen graph
minus six edges). In particular $\Nanti$ is odd.
\end{proposition}

\noindent\textbf{Status.}  The equivalence ``no intrinsic-$q$
$\iff\ell|_R\not\equiv0$'' is the proven Lemma~\ref{lem:skeleton}(iv). The value
$\Nanti=9$ is a \emph{realizability} statement: it holds with $0$ exceptions over
$154$ no-$q$ configurations across two disjoint seed families ($2695$
$(8,6,2)$-configurations in total; \texttt{noq\_odd\_proof.py}), but it is not a
formal consequence of the quadratic form alone. Indeed
Lemma~\ref{lem:skeleton}(v) gives $\Nanti=o(\mathbf 1)$, and in the intrinsic-$q$
\emph{existing} bucket $o(\mathbf 1)$ realises both parities (the
$404/103$ split of Proposition~\ref{prop:intrinsicq}); so $o$ and $\ell$ do not by
themselves pin the value. Proposition~\ref{prop:noq_odd} is therefore of the same
character as the realizability statement $\rank G=4\Rightarrow\Nanti=6$
(Remark~\ref{rem:lowbound})---a constraint on which anti-graphs the Lagrangian
geometry actually produces at fixed symplectic rank.

\begin{remark}
The skeleton localizes exactly what an eventual proof must supply. By
Lemma~\ref{lem:skeleton} the no-$q$ hypothesis is $\ell|_R\not\equiv0$, i.e.\ a ray
dependency $\sum_{i\in a}r_i=0$ with an odd internal anticommuting count
$o(a)=1$; one must show this forces the global count $o(\mathbf 1)$ to the specific
value $9$ (equivalently, the anti-graph to the degree class above). This is a
combinatorial-geometric realizability question about proper $K_5$, not an identity
of the quadratic form $o$.
\end{remark}

\begin{remark}[A reduction frame, and why the realizability layer is itself stratified]
\label{rem:reduction}
Proposition~\ref{prop:noq_odd} and the lower bound $\rank G=4\Rightarrow\Nanti=6$
(Remark~\ref{rem:lowbound}) are two instances of one reduction. The radical
quotient $\bar W=W/\rad(\omega|_W)$ is a \emph{non-degenerate} symplectic space of
dimension $\rank G$, and each Lagrangian $L_a$ descends to an isotropic subspace
$\bar L_a\subseteq\bar W$ (its rays lie in the isotropic $L_a$, hence commute).
A proper $K_5$ thus projects to a five-Lagrangian configuration in
$\Sp(\rank G,\mathbb F_2)$, and $\Nanti$ is the anticommuting count of the
projected rays. The two problems become reduced-configuration realizability
questions: $\rank G=4$ lands in $\Sp(4,\mathbb F_2)$, where $\Nanti=6$ is forced
(Theorem~\ref{thm:rank4}, Section~\ref{sec:rank4}), while $\rank G=6$ lands
in $\Sp(6,\mathbb F_2)$ (where the value is forced only after the extra selector
$\ell|_R\not\equiv0$).

Crucially the forcing is itself \emph{dimension-stratified}. Empirically $\Nanti$
is pinned by low-arity data---the stratum $(\dim W,\rank G,\dim\rad W)$, the
projected-Lagrangian dimensions, and the bit $\ell|_R$---for every stratum through
$\dim W=8$, but by \emph{none} of these invariants at $\dim W=9$ (there it splits
across every refinement we tried: stratum $\times$ projected-dimensions $\times
\ell|_R$). So the two realizability statements are of the same kind but do
\textbf{not} merge into a single dimension-free theorem: the realizability layer
softens with dimension exactly as the rigidity does---a structural echo, one level
down, of the uniqueness of $n=4$. This reinforces the modulus thesis: not only is
the fiber unclassified by arity-$\le4$ relative-position data
(Theorem~\ref{thm:modulus}), even the \emph{coarse value} $\Nanti$ ceases to be
determined by the reduced geometry once $\dim W\ge9$.
\end{remark}

% ------------------------------------------------------------------
\section{Partial Rigidity Results}\label{sec:partial}
% ------------------------------------------------------------------

Despite the modulus obstruction, several deterministic implications survive.

\begin{conjecture}[Threshold]\label{conj:threshold}
For any proper $K_5$ in $\Sp(10,\mathbb F_2)$: if $n_{\mathrm{odd}}\le5$ then
$\Nanti$ is even.
\end{conjecture}

\textbf{Empirical support.}  Zero violations across 455 proper $K_5$ in the
$n_{\mathrm{odd}}\le5$ range (8 independent seeds, seed family $2025+101s$).
The threshold is sharp: $n_{\mathrm{odd}}=6$ produces both parities
(empirically 230 even / 7 odd; the 7 odd cases are not concentrated in a single
hypergraph type).

\begin{remark}
A proof of Conjecture~\ref{conj:threshold} would require relating the odd-triple
hypergraph to $\Nanti\bmod2$ via a global identity.  No clean modular formula
$\sum\mu_{abc}\equiv\Nanti\bmod2$ has been found; the implication appears
one-directional (a sufficient condition for evenness, not an equivalence).
\end{remark}

% ------------------------------------------------------------------
\section{The rank-4 stratum: a descent to the doily}\label{sec:rank4}
% ------------------------------------------------------------------

In the lowest non-trivial rank the $n=4$ rigidity \emph{survives}, and we can pin it
exactly. The mechanism is a descent: the whole proper $K_5$ projects along the radical
quotient into a small symplectic space, and at $\rank G=4$ that space is a classical
object---the $S_6$ doily---in which the question ``$\Nanti=6$?'' becomes finite and
yields to a short counting argument. The result is the first $n\ge5$ realizability
statement we can prove outright.

First, $\rank G=4$ bounds the radical. Writing $M$ for the $10\times10$ matrix of rays
and $\Omega$ for the form, $G=M\Omega M^{\top}$, so by the Sylvester rank inequality
\[
  \rank G \;=\; \rank(M\Omega M^{\top}) \;\ge\; 2\,\rank M-10 \;=\; 2\dim W-10 ,
\]
whence $\dim W\le 7$ and $\dim\rad(W)=\dim W-\rank G\le 3$. (Equality $\dim W=7$ holds
in every sampled configuration, but only the bound $\dim\rad(W)\le3$ is used below.)

The reduction (Remark~\ref{rem:reduction}) sends the configuration to
$\bar W=W/\rad(W)$, a non-degenerate symplectic space of dimension $\rank G=4$, i.e.\
$\Sp(4,\mathbb F_2)\cong S_6$---the syntheme--duad doily $GQ(2,2)$ (points $=$ duads,
i.e.\ $2$-subsets of $\{1,\dots,6\}$; lines $=$ synthemes, i.e.\ perfect matchings;
$\omega(D,D')=|D\cap D'|\bmod2$). The five Lagrangians become five synthemes; the ray
$r_{ab}$ projects to the duad shared by $S_a,S_b$, or to $0$ if they are edge-disjoint
(then $v_{ab}\in\rad(W)$). An exhaustive sweep of all five-syntheme configurations
satisfying the projection axioms (\texttt{doily\_rank4.py}) gives $\Nanti\in\{5,6\}$,
with $\Nanti=5$ \emph{iff} the meeting-graph on the five indices is the pentagon
$C_5$---equivalently the five non-meeting pairs form the complementary pentagram, whose
rays lie in $\rad(W)$ (dimension $\le3$). Two lemmas then exclude that case.

\begin{lemma}[Three-term relations are Fano stars]\label{lem:3term}
Among any three distinct rays of a proper $K_5$, the only $\mathbb F_2$-linear
dependency $v_p+v_q+v_r=0$ is a \emph{Fano star}: the three rays share a common
Lagrangian, $v_{ij}+v_{ik}+v_{il}=0$.
\end{lemma}

\begin{proof}
Classify by the index pattern of $\{p,q,r\}$. If the three pairs share a common
index ($\{ij,ik,il\}$) it is the stated star. Otherwise some two of them, say
$v_p,v_q$, share an index $y$ while the third $v_r=\{u,v\}$ does not equal a known
ray: then $v_r=v_p+v_q\in L_y$, so $v_r\in L_y\cap L_u=\langle v_{yu}\rangle$, forcing
$v_r=v_{yu}$, against distinctness. (Concretely: triangle $\{xy,xz,yz\}$, path
$\{xy,yz,zw\}$, and the disjoint pattern $\{xy,yz,uv\}$ all fall here.) Hence only the
star survives. Fano stars do occur for $n\ge5$, where B0 fails ($\approx3.8\%$ of
configurations carry a Fano vertex).
\end{proof}

\begin{proposition}[The $C_5$ pentagon does not lift]\label{prop:c5}
No proper $K_5$ in $\Sp(10,\mathbb F_2)$ has five rays whose index-pairs form a
pentagram $\{(i,i+2):i\in\mathbb Z_5\}$ all contained in a subspace of dimension
$\le3$. Consequently the $C_5$ doily configuration is not realizable.
\end{proposition}

\begin{proof}
The idea is a pigeonhole on relation-weights: five rays squeezed into three dimensions
must satisfy many relations, but the pentagram is too ``spread out'' (no three rays
share a Lagrangian) to admit any short one---and short relations are exactly what a
low-dimensional crowd of vectors is forced to produce.

Suppose the five pentagram rays $v_{i,i+2}$ lie in a subspace of dimension $\le3$ (in
our application, $\rad(W)$, of dimension $\le3$ by the Sylvester bound). Being five
vectors in dimension $\le3$, their space of linear relations has dimension $\ge2$. Each index occurs in exactly two
pentagram pairs, so no three of the five rays share a common index; by
Lemma~\ref{lem:3term} there is no Fano star among them, hence no relation of weight
$\le3$ (weight $1$ would be a zero ray, weight $2$ two equal rays). Thus every
nonzero relation has weight $\ge4$ on the five coordinates. But a $\ge2$-dimensional
space contains two independent weight-$\ge4$ vectors $r_1,r_2\in\mathbb F_2^5$; since
$|\mathrm{supp}\,r_1|+|\mathrm{supp}\,r_2|\ge8>5$ we have
$|\mathrm{supp}\,r_1\cap\mathrm{supp}\,r_2|\ge3$, so $r_1+r_2$ has weight $\le2$---a
nonzero relation of weight $\le2$, contradiction.
\end{proof}

\begin{theorem}[Rank-4 value]\label{thm:rank4}
Every proper $K_5$ in $\Sp(10,\mathbb F_2)$ with $\rank G=4$ satisfies $\Nanti=6$.
\end{theorem}

\begin{proof}
By the Sylvester bound above, $\rank G=4$ gives $\dim\rad(W)\le3$, and the doily
descent sends the configuration into $\Sp(4,\mathbb F_2)$. The exhaustive
classification (\texttt{doily\_rank4.py}) leaves only $\Nanti\in\{5,6\}$, with
$\Nanti=5$ \emph{iff} the five non-meeting pairs form a pentagram whose rays lie in
$\rad(W)$. Since $\dim\rad(W)\le3$, Proposition~\ref{prop:c5} forbids exactly that, so
$\Nanti=6$. (Independently verified $176/176$ across three disjoint seed families.)
\end{proof}

\begin{remark}
The doily sweep ranges over configurations of five \emph{distinct} synthemes. In the
measure-small degenerate case where two Lagrangians project to the same line of
$\bar W$ the same radical-counting applies (the non-meeting rays still lie in
$\rad(W)$); we have additionally verified $\Nanti=6$ directly there ($0$ exceptions
over all sampled $(7,4,3)$-configurations).
\end{remark}

\begin{remark}\label{rem:lowbound}
Theorem~\ref{thm:rank4} is genuinely a \emph{realizability} statement: the matching
lower bound $\Nanti\ge6$ is not graph-theoretic (rank-4 Petersen subgraphs with
$2$--$5$ edges abound---$987$ of $1027$), and it is the Lagrangian geometry, distilled
into the radical-counting of Proposition~\ref{prop:c5}, that forces the maximum. The
two realised anti-graphs are the $6$-cycle $C_6$ ($\approx91\%$) and a $7$-vertex tree
($\approx9\%$), both with $6$ edges; which of the two occurs is the only part of the
rank-4 picture left open.
\end{remark}

\begin{proposition}[Adjacent low-rank strata]\label{prop:lowrank}
$\rank G=0\Leftrightarrow\Nanti=0$ (tautological); and $\rank G=2\Rightarrow
\Nanti\in\{1,3\}$ (odd; $\le3$ rigorous by Theorem~\ref{thm:upbound}, the values
$\{1,3\}$ observed $8/8$, proof open).
\end{proposition}

% ------------------------------------------------------------------
\section{Open Problems}\label{sec:open}
% ------------------------------------------------------------------

\begin{enumerate}
\item \textbf{Prove Conjecture~\ref{conj:threshold}} ($n_{\mathrm{odd}}\le5\Rightarrow
  \Nanti$ even).  A global identity relating the odd-triple hypergraph to $\Nanti\bmod2$
  would suffice.

\item \textbf{Finer rank-4 anti-graph type ($C_6$ vs.\ tree).}
  The rank-4 \emph{value} is now settled: $\rank G=4\Rightarrow\Nanti=6$
  (Theorem~\ref{thm:rank4}, via the $S_6$-doily reduction and the
  non-lifting of the $C_5$ pentagon, Proposition~\ref{prop:c5}). What remains is the
  finer question of \emph{which} $6$-edge anti-graph occurs: a $6$-cycle $C_6$
  (na-vector $(2,2,2,0,0)$, $\approx91\%$) or a $7$-vertex tree
  ($(2,2,1,1,0)$, $\approx9\%$). A clean invariant separating the two would complete
  the rank-4 picture.

\item \textbf{Complete fiber classification or prove it is a genuine modulus.}
  Theorem~\ref{thm:modulus} shows arity $\le4$ is insufficient.  The question remains
  whether a Wall-type invariant at arity 5 (involving all five Lagrangians) classifies
  the na-vector, or whether the fiber is truly a modulus (no finite-arity classifier).
  The quad-bit saturation ($99.5\%$ at value 1) suggests the natural arity-4 object is
  degenerate; a \emph{different} arity-4 invariant might behave better.

\item \textbf{Prove the value in Proposition~\ref{prop:noq_odd}}
  (no intrinsic-$q$ $\Rightarrow\Nanti=9$ in the $(8,6,2)$-stratum). By
  Lemma~\ref{lem:skeleton} the hypothesis reduces to a ray dependency with odd
  internal anticommuting count; what is open is the realizability step forcing the
  global anti-graph to the degree class $(0,0,2,2,2,2,2,2,3,3)$. Via
  Remark~\ref{rem:reduction} this is the $\Sp(6,\mathbb F_2)$ instance of the
  reduction (item~2 is the $\Sp(4,\mathbb F_2)$ instance). The two are the same
  \emph{kind} of question but---as the dimension-stratification in
  Remark~\ref{rem:reduction} shows---do not collapse to one dimension-free theorem:
  the $\Sp(6)$ case additionally requires the selector $\ell|_R\not\equiv0$, and the
  forcing fails entirely by $\dim W=9$. A proof should treat each reduced dimension
  on its own terms.

\item \textbf{$n=6$ and beyond.}  Does the order parameter $\rad(\omega|_W)$ persist
  as an attractor at $n=6$?  Is there a second structural transition?  Is there an
  infinite descent (ever weaker rigidity as $n\to\infty$) or does the structure
  stabilise?

\item \textbf{Algebraic proof of Key Lemma from Paper~XVIII} (rank$(G)\in\{0,8\}
  \Leftrightarrow\Sigma_m=0$).  Verified computationally for all 840 proper $K_5$
  up to symmetry.  An algebraic proof would clarify whether this is a $n=4$-specific
  coincidence or the shadow of a deeper identity.
\end{enumerate}

% ------------------------------------------------------------------
\section{Connections to Previous Papers}\label{sec:connections}
% ------------------------------------------------------------------

\textbf{Paper XI}~\cite{PaperXI} (Quadratic refinement) introduced the Arf invariant as a potential
classifier of parity.  Section~\ref{sec:arf} shows that for $n\ge5$ the intrinsic-$q$
is generically well-defined but its Arf is non-discriminating.  The Arf machinery is
powerful for $n=4$ (where it gives an alternative characterization of the Key Lemma)
but cannot handle the $n\ge5$ fiber.

\textbf{Paper XII}~\cite{PaperXII} ($T$-vector) introduced $T=\bigoplus v_{ab}$.
In the $n=4$ setting the T-vector is the primary obstruction
(\cite{PaperXVIII}, Section~3).  For $n\ge5$ the T-vector is no longer universally
zero; its relation to $\rad(W)$ and to the odd-triple count $n_{\mathrm{odd}}$ is
an open direction.

\textbf{Paper XIII}~\cite{PaperXIII} (Maslov index) introduced the Maslov triple index
as the relative-position invariant of three Lagrangians.  The present paper confirms
this is the correct level of analysis for $n\ge5$---the Maslov triple provides real
partial rigidity ($n_{\mathrm{odd}}\le5\Rightarrow\Nanti$ even) but is insufficient
for complete classification.  The Wall-type extension (arity 4 and beyond) is the
natural continuation.

\textbf{Paper XVII}~\cite{PaperXVII} established that for an all-Fano proper $K_5$ in
$\Sp(6,\mathbb F_2)$ the anticommutation graph is the full Petersen graph $K(5,2)$
($\Nanti=15$), and proved the B0 lemma.  This is the all-Fano half of the $n=3$
trichotomy vertex in Theorem~\ref{thm:trichotomy}; the generic $n=3$ proper $K_5$
(outside the all-Fano $\approx16\%$) instead has $\Nanti=12$.

\textbf{Paper XVIII}~\cite{PaperXVIII} proved $\Nanti=10$ universally for $n=4$, and
identified the transversal property (B0+B1+B2) as the correct $K_4$-level proof.
The present paper shows the B2 failure for $n\ge5$ is precise ($W^\perp$ has 15
nonzero vectors, well above the pigeonhole threshold of 3), and that B1 fails because
the commuting-pairing span has dimension exactly one below the Lagrangian dimension.

\subsection*{Synthesis}

The two halves of this paper answer the opening question with a single picture. The
$n=4$ rigidity does not vanish at $n\ge5$; it \emph{stratifies}. Generically---in the
high-rank strata---it dissolves into a modulus that no arity-$\le4$ invariant resolves
(Theorem~\ref{thm:modulus}). In the low-rank strata it persists and is provable, via
the symplectic descent $W\rightsquigarrow\bar W=\Sp(\rank G,\mathbb F_2)$: at
$\rank G=4$ the descent reaches the $S_6$ doily and pins $\Nanti=6$
(Theorem~\ref{thm:rank4}). The same dimension that makes the descent land in a
\emph{small} space when $\rank G$ is small is what lets the modulus spread when
$\rank G$ is large; and the very Fano stars that the rank-4 proof must avoid are the
debris of the B0-collapse that ended $n=4$. Rigidity, modulus, and the
Kochen--Specker obstruction of Papers~XI--XVIII thus appear as facets of one
dimension-graded structure.

Finally, we flag a forward connection. Paper~IV~\cite{PaperIV} conjectured, from the
cohomological ladder $H^1\to H^2\to H^3$, that a genuinely third-order
(``Borromean'', non-associative) obstruction can appear only for $N\ge5$ qubits---and
derived this threshold from a \emph{global symplectic capacity bound} that collapses
the $5$-MASA nerve $\partial\Delta^4$ at $N\le4$. That nerve is precisely a proper
$K_5$, its collapse is precisely the B0/B1/B2 rigidity of Paper~XVIII, and its release
at $n=5$ is the regime of the present paper. We therefore conjecture that the modulus
of Theorem~\ref{thm:modulus}---unresolved by any arity-$\le4$ (degree-$\le3$
relative-position) invariant---is the symplectic realisation of that $H^3$ class, with
the quadruple invariant $q_4$ playing the role of the degree-3 cochain. Making this
identification precise is the natural sequel.

% ------------------------------------------------------------------
\section{Rigor Self-Assessment}\label{sec:rigor}
% ------------------------------------------------------------------

\begin{center}
\begin{tabular}{p{5cm}p{3.5cm}p{5cm}}
\toprule
Claim & Status & Note \\
\midrule
Trichotomy (B1/B2 dimension arithmetic) & Rigorous & Algebraic; cited from XVIII \\
Upper-bound ladder & Rigorous & Exhaustive $2^{15}$ enumeration \\
Lower bound $\rank G=4\Rightarrow\Nanti=6$ & \textbf{Rigorous} & Doily reduction + radical-counting (Prop.~\ref{prop:c5}); $C_5$ pentagon excluded \\
$C_5$ pentagon non-liftable (Prop.~\ref{prop:c5}) & Rigorous & Counting in $\rad(W)$; Lemma~\ref{lem:3term} (3-term $=$ Fano star) \\
Modulus witness & Rigorous (computational) & Explicit pair; independently reproducible \\
$n_{\mathrm{odd}}\le5\Rightarrow\Nanti$ even & Conjectural (455/455) & Proof open \\
frame-$q$ non-invariance & Rigorous & Algebraic + 900-case verification \\
Structural skeleton, Lemma~\ref{lem:skeleton} (i)--(v) & Rigorous & Elementary $\mathbb F_2$ linear algebra; coisotropic reduction + T-vector \\
Intrinsic-$q$ exists $\iff\ell|_R\equiv0$ & Rigorous & Lemma~\ref{lem:skeleton}(iv); $0/2695$ vs.\ direct solver \\
Arf non-discriminating in $(8,6,2)$ & Rigorous (computational) & 404+103 split \\
No intrinsic-$q\Rightarrow\Nanti=9$ (hence odd) & Conjectural (154/154) & Realizability; reduction proven (Lemma~\ref{lem:skeleton}), value open \\
$\mu$ symmetric in $S_3$ & Rigorous (computational) & 25-config self-check \\
Order parameter $\rad(W)$ & Descriptive & Monotone, not a classifier \\
\bottomrule
\end{tabular}
\end{center}

All computational results are reproducible via the supplementary
scripts~\cite{PaperXIXsuppl}:
\begin{verbatim}
cd supplementary/paper19
python3 lowrank_rigidity.py 5 3000 600 8    # ladder + strata
python3 maslov_probe.py     5 2500 350 8    # witness, mu, purity
python3 quad_refine.py      5 2500 350 6    # Arf ruling-out
python3 noq_odd_proof.py    5 4000 600 8    # skeleton + no-q=>N_anti=9
python3 reduction_frame.py  5 4000 500 10   # reduction + forced strata
python3 doily_rank4.py                      # rank-4 via S6 doily
python3 rank4_lemmas.py     5 4000 500 10   # rank-4 lemmas (I),(II) + chain
python3 rank4_lemmaII.py    5 4000 400 8    # (II) anatomy (Fano star)
\end{verbatim}

\bigskip
\noindent\textbf{Acknowledgements.}
Computations performed with \texttt{co-nlang} Research infrastructure.

\begin{thebibliography}{99}

\bibitem[PaperIV]{PaperIV}
Z.-L. Chen,
``Lyndon--Hochschild--Serre Transgression and the $H^3$ Frontier:
From Group Cohomology to Borromean Contextuality (Paper~IV),''
\textit{Zenodo, DOI: 10.5281/zenodo.20482595} (2026).

\bibitem[PaperXI]{PaperXI}
Z.-L. Chen,
``Quadratic Refinement and the Parity of Mermin Pentagrams:
The $\beta$-Formula, Geometric Decomposition, and the
Kochen--Specker Obstruction from Equiangular Geometry,''
\textit{Zenodo, DOI: 10.5281/zenodo.20482595} (2026).

\bibitem[PaperXII]{PaperXII}
Z.-L. Chen,
``The T-Vector Theorem: Symplectic Characteristic Elements
and the Kochen--Specker Obstruction,''
\textit{Zenodo, DOI: 10.5281/zenodo.20490118} (2026).

\bibitem[PaperXIII]{PaperXIII}
Z.-L. Chen,
``The Maslov Index and the Kochen--Specker Obstruction:
Negative Results, the $k$-Profile Theorem, and Orbit~4 Anomaly,''
\textit{Zenodo, DOI: 10.5281/zenodo.20496513} (2026).

\bibitem[PaperXIV]{PaperXIV}
Z.-L. Chen,
``Stabilizer Algebra and the $k$-Profile Theorem for Mermin Pentagrams,''
\textit{Zenodo, DOI: 10.5281/zenodo.20502868} (2026).

\bibitem[PaperXV]{PaperXV}
Z.-L. Chen,
``The Weil Representation and the $S_3$ Lifting Classification
for Mermin Pentagrams,''
\textit{Zenodo, DOI: 10.5281/zenodo.20519654} (2026).

\bibitem[PaperXVI]{PaperXVI}
Z.-L. Chen,
``The Weyl Product Identity and the Algebraic Origin
of the Kochen--Specker Obstruction,''
\textit{Zenodo, DOI: 10.5281/zenodo.20519733} (2026).

\bibitem[PaperXVII]{PaperXVII}
Z.-L. Chen,
``The Cross-Context Anticommutation Theorem
and the Algebraic Origin of the Kochen--Specker Obstruction,''
\textit{Zenodo, DOI: 10.5281/zenodo.20530757} (2026).

\bibitem[PaperXVIII]{PaperXVIII}
Z.-L. Chen,
``Mermin Pentagrams in $\Sp(8,\mathbb{F}_2)$:
Universal $N_{\mathrm{anti}}=10$ and the Gram Rank Selection Principle,''
\textit{Zenodo, DOI: 10.5281/zenodo.20579239} (2026).

\bibitem[PaperXIXsuppl]{PaperXIXsuppl}
Z.-L. Chen,
``Supplementary computational scripts for Paper~XIX,''
\textit{Zenodo, DOI: 10.5281/zenodo.20579390} (2026).

\end{thebibliography}

\end{document}
